% Source map assembled before drafting the chapters, 7 October 2026.
\chapter{Source map and scope of the reconstruction}\label{app:sources}

This map was assembled before the chapters were drafted. References were
checked in the passages listed below, including the local Fantechi
transcription and the local Behrend--Manin text. Page numbers refer to the
linked version unless explicitly called journal pages. This is a map of
what each source actually supplies, not a claim that one reference proves
every step of these notes.

\begin{description}[style=nextline]
\item[Categories, pullbacks, and the two stack conditions.]
Fantechi, \emph{Stacks for Everybody} \cite{Fantechi}, Sections 3.2,
3.5, 4.1--4.3 and 6.3: Definition 3.1 (source PDF p.~5),
Definition 4.1 (p.~6), Definition 4.2 (p.~6), Definition 4.3 (p.~7).
The local file is
\path{../../Papers/Stacks/StacksForEverybody.tex}, relative to this
directory. Its labels are \texttt{def:3.1}, \texttt{def:4.1},
\texttt{def:4.2}, and \texttt{def:4.3}.
This is a later pedagogical exposition using bundles; our curves and maps
are applications, not examples whose full proofs appear in that article.

\item[Gluing morphisms, and hence isomorphisms.]
Vistoli, \emph{Notes on Grothendieck topologies, fibered categories and
descent theory} \cite{Vistoli}, Section 4.3, Proposition 4.31,
pp.~91--92, and the representable-sheaf theorem used there;
Stacks Project \cite{Stacks}, \MCtag{023Q}, Lemma 35.13.7.
The latter checks the exact statement: scheme-valued representable
functors are fpqc sheaves. We apply it on a cover of the \emph{source
total space}, glue the inverse separately, and check markings and target
maps by equality after a cover. No representability theorem for the
isomorphism functor is needed to prove this stack condition.

\item[Effective descent of the total spaces.]
SGA~1 \cite{SGA1}, Expos\'e VIII, Section 7, Proposition 7.8, is the
scheme-with-ample-invertible-sheaf descent theorem. Vistoli,
Section 4.3.3, Theorem 4.38 and its proof, pp.~96--103, gives a detailed
later explanation using a functorial relatively ample invertible sheaf.
His Example 4.39 treats smooth curves, not our entire pointed nodal/map
problem. Our reconstruction supplies the extra checks: the functorial
bundles $\omega_{C/S}(\sum\sigma_a)$ and
$\omega_{C/S}(\sum\sigma_a)\otimes f^*A^{\otimes3}$, their cocycles,
and descent of the sections and of $f$.
Vistoli, Proposition 1.15, pp.~11--12, records descent of the required
morphism properties. Stacks Project \MCtag{0E6R}, Lemma 109.19.3,
checks the invertibility and arbitrary-base-change property of the
relative dualizing sheaf for Gorenstein curve families.

\item[Stable curves, embeddings, and geometric properties.]
Deligne--Mumford \cite{DM}, Section 1, Definition (1.1),
Theorem (1.2) and its corollary, and Lemma (1.4), journal pp.~76--81,
give the unpointed stability/embedding/automorphism ingredients;
Section 5, Proposition (5.1) and Theorem (5.2), pp.~104--105, give
the unpointed stack and its properness and smoothness.
Knudsen \cite{Knudsen}, Theorem 2.7, journal p.~179, is the pointed
proper smooth Deligne--Mumford theorem.
These are original research sources whose indicated statements were
checked. Alper \cite{Alper}, in the dated 5 January 2026 version,
Propositions 6.3.23--6.3.28, pp.~231--235,
Theorems 6.4.6 and 6.4.16, pp.~238--241, and
Section 6.5, pp.~242--252, supplies the later detailed exposition.
Our atlas proof applies his embedding-lifting argument directly
to stable pointed families. Edidin \cite{Edidin}, Section 3,
offers a supplementary unpointed Hilbert-scheme exposition.
The guided proofs explicitly retain Hilbert
representability, bounded embeddings, deformation theory of nodal
curves, and stable reduction as external theorems. They are not
deductions from the two stack axioms alone.

\item[Stable maps: definitions and bounded parameter spaces.]
Fulton--Pandharipande \cite{FP}, Section 1.1, pp.~10--11, gives
families, isomorphisms, and the ample stability bundle;
Section 2.3, pp.~13--14, gives bounded embeddings and the Hilbert
incidence construction. Section 5.1, Lemma 8, pp.~26--27, handles
factorization through the target. This is an exposition and primarily
a construction of \emph{coarse spaces}; Section 1.2 explicitly sets
aside the stack viewpoint. Identifying the corresponding quotient
\emph{stack}, and verifying Fantechi's axioms, are our guided
reconstruction with the descent references above.
Behrend--Manin \cite{BM}, Definition 2.1, p.~12, supports the
degree-functional convention for $\beta$; Definitions 3.2 and 3.13,
pp.~13--14 and 25, give stable-map and stack notation.
Definition 3.13 asserts the stack property without supplying the
step-by-step descent proof requested here.

\item[Stable maps: Deligne--Mumford, finite type, and properness.]
Behrend--Manin, Theorem 3.14, p.~25, Proposition 4.1 and
Lemma 4.2, pp.~26--27, and Corollary 4.8, p.~28, support these
properties for the stated projective target and characteristic
assumptions. This is an early research treatment; no absolute
priority claim is needed here. Its properness argument invokes
earlier preliminary notes. Fulton--Pandharipande, Section 4.2,
Propositions 5 and 6, pp.~19--21, supplies readable proofs about
extending isomorphisms and families, respectively. The latter
uses semistable reduction as an external theorem.
Our properness argument passes from those family statements to
the stack valuative criterion. Smoothness is established in the
explicit degree-one plane example by a direct identification with
the dual projective plane; no general smoothness assertion is made.
\end{description}

\section{How the search was evaluated}

The local search included all of \path{DoctoralNotebook/Papers/}:
in particular the Alper PDF, Fantechi transcription, Behrend--Manin,
Behrend's Gromov--Witten paper, the Faber--Pandharipande text, and the
later virtual-class and pseudostability material. The last groups
were not used as substitutes for foundations or brought into the
present scope. The earlier $BG$ and Riemann--Roch summaries supplied
the memoir/header conventions, checkpoints, and exercise format.

The public search used the suggested authors as leads and also found
Vistoli, Edidin, Schmitt, and student seminar notes. It distinguished
foundational construction from later work that uses these stacks.
Supplementary sources and their limitations are recorded in
Section~\ref{sec:reading-audit}. No advisor--student relationship is
asserted. When a proof below invokes a substantial theorem, it is
labelled as an external input rather than concealed in a phrase such
as ``the objects glue.''
